\documentclass[11pt]{article}
\usepackage[T1]{fontenc}
\usepackage{lmodern,amsmath,amssymb,amsthm,mathtools}
\usepackage[margin=1.05in]{geometry}
\usepackage[hidelinks]{hyperref}
\newtheorem{theorem}{Theorem}[section]
\newtheorem{lemma}[theorem]{Lemma}
\newtheorem{proposition}[theorem]{Proposition}
\newtheorem{corollary}[theorem]{Corollary}
\newcommand{\C}{\mathbb C}\newcommand{\R}{\mathbb R}
\newcommand{\E}{\mathbb E}\newcommand{\Pp}{\mathbb P}
\newcommand{\code}[1]{{\small\nolinkurl{#1}}}
\title{Complex fourth-root approximations:\\analytic limits, compact escape and a fixed variation test}
\author{Project Stardust}
\date{Reviewed editorial draft, October 29, 2026}
\begin{document}
\maketitle
\begin{abstract}
Uniform complex fourth-root steps yield genuine positive probability laws
whose holomorphic exponential transforms approach the biharmonic heat
symbol. Their absolute second moments diverge, and exact fourth moments
give a uniform positive mass outside each compact set, eventually along
every refinement. Thus the transform limit is compatible with the absence
of any weak probability-law limit. Poissonization preserves both features.
Normalized complex exponential tilts have mass one and convergent
holomorphic transforms, but their total variations diverge. We exhibit a
single bounded continuous test, independent of the approximation index,
whose integral equals the total variation exactly. The paper reconstructs
verified fourth-order results with elementary proofs and explicit
constants. The analytic-versus-probabilistic distinction is established
in prior work; this is a formalization companion, and no priority claim
is made.
\end{abstract}

\section{Scope and main conclusions}
Let $\varepsilon_j,\delta_j$ be mutually independent symmetric signs and set
\[
 \zeta_j=\varepsilon_j+i\delta_j,\qquad
 S_{n,a}=a\sum_{j=1}^n\zeta_j,\qquad
 a_n(T)=(6T/n)^{1/4}\quad(n\geq1).
\]
The four values of $\zeta_j$ are the fourth roots of $-4$.
Write $\mu_n^T$ for the law of $S_{n,a_n(T)}$ on $\C$ with its usual
Euclidean topology. For $r\geq0$, let $N_r$ be Poisson with mean $r$,
independent of the signs, and let $\mu_{r,a}$ be the law of
\[
 C_{r,a}=a\sum_{j=1}^{N_r}\zeta_j.
\]
All these laws are ordinary positive probability measures.

\begin{theorem}\label{thm:positive}
For $T>0$ and each real $\eta$,
\[
 \int e^{i\eta z}\,d\mu_n^T(z)\longrightarrow e^{-T\eta^4},\qquad
 \int e^{i\eta z}\,d\mu_{n,a_n(T)}(z)\longrightarrow e^{-T\eta^4}.
\]
Nevertheless, for every compact $K\subset\C$, eventually in $n$,
\[
 \mu_n^T(K^c)\geq\tfrac18,\qquad
 \mu_{n,a_n(T)}(K^c)\geq\tfrac1{12}.
\]
Neither family is tight. Neither has a weak probability-law limit along
any reindexing $n=\varphi(k)$ with $\varphi(k)\to\infty$.
\end{theorem}

The transforms in this statement test $e^{i\eta z}$, which is unbounded on
$\C$ when $\eta\ne0$. They are not the characteristic functions
$e^{i(u\operatorname{Re}z+v\operatorname{Im}z)}$ used in weak convergence
on $\R^2$. This domain distinction is essential.

For the tilted result use the second, rate-scaled normalization
$b_n=a_n(1)$, $r_n=nT$. Define real functions
\[
 F(x)=\cos x\cosh x,\qquad \psi_a(\xi)=F(a\xi)-1,
\]
and the finite complex measure
\begin{equation}\label{eq:tilt}
 d\nu_n(z)=e^{-r_n\psi_{b_n}(\xi)}e^{i\xi z}\,d\mu_{r_n,b_n}(z).
\end{equation}
Here and below total variation means the mass $|\nu|(\C)$, with no factor
$1/2$.

\begin{theorem}\label{thm:tilt}
Fix $T>0$ and $\xi\in\R\setminus\{0\}$. Each $\nu_n$ has mass one, and
for every real $\eta$,
\begin{equation}\label{eq:tiltlimit}
 \int e^{i\eta z}\,d\nu_n(z)
 \longrightarrow e^{-T((\eta+\xi)^4-\xi^4)}.
\end{equation}
The fixed function $f_\xi(z)=e^{-i\xi\operatorname{Re}z}$ is continuous
and has modulus one, and satisfies
\begin{equation}\label{eq:fixed}
 \int f_\xi\,d\nu_n=|\nu_n|(\C)
 =\exp\!\left(r_n\cosh(b_n\xi)(1-\cos(b_n\xi))\right)
 \longrightarrow+\infty.
\end{equation}
Consequently these complex measures do not converge against all bounded
continuous tests to any finite complex measure, even along a divergent
reindexing.
\end{theorem}

\paragraph{Relationship with prior work.}
Bonaccorsi and Mazzucchi~\cite{BM} construct higher-order analytic limits
from complex root walks. Their Theorem 2 treats exponential transforms;
their Remark 1 discusses escape from bounded regions, and their discussion
before Theorem 7 distinguishes the analytic limit from an ordinary limit
process. Our rotated fourth roots and normalization are a fourth-order
instance of that construction. We do not assert that their analytic
framework supplies an ordinary weak limit. The contribution of this
companion is a complete reconstruction of the frozen formal results:
finite moment identities, explicit escape bounds, Poissonized statements,
and a fixed bounded test for actual complex coordinate measures. These
are meaningful precise distinctions, but their independent novelty has
not been established. The later analytic It\^o framework~\cite{BCM} is
also separate from weak convergence of the present probability laws.

\section{Transforms and integrability}
\begin{lemma}\label{lem:step}
For real $a,\eta$,
\begin{align}
 \E e^{i\eta a\zeta_1}&=F(a\eta),\label{eq:step}\\
 \E\left|e^{i\eta a\zeta_1}\right|&=\cosh(a\eta).
 \label{eq:stepnorm}
\end{align}
Consequently
\begin{align}
 \E e^{i\eta S_{n,a}}&=F(a\eta)^n,\label{eq:walktransform}\\
 \E e^{i\eta C_{r,a}}&=\exp(r(F(a\eta)-1)),\label{eq:compound}\\
 \E\left|e^{i\eta C_{r,a}}\right|&=\exp(r(\cosh(a\eta)-1)).
 \label{eq:compoundnorm}
\end{align}
All integrals in these formulas are absolutely integrable.
\end{lemma}
\begin{proof}
Independence of the two signs gives
$\E e^{ix(\varepsilon+i\delta)}=(\E e^{ix\varepsilon})
(\E e^{-x\delta})=\cos x\cosh x$. The modulus is $e^{-x\delta}$,
which gives \eqref{eq:stepnorm}. Independence over $j$ gives products
for the walk, including $\E|e^{i\eta S_{n,a}}|=\cosh(a\eta)^n$.
Conditioning on $N_r=m$ and summing the Poisson series gives
\[
 e^{-r}\sum_{m=0}^\infty\frac{r^m}{m!}F(a\eta)^m
 =e^{r(F(a\eta)-1)}.
\]
Absolute convergence follows either from the same series with
$|F(a\eta)|$ or from the nonnegative modulus series
$e^{-r}\sum_m r^m\cosh(a\eta)^m/m!$.
This also proves \eqref{eq:compoundnorm} and justifies conditioning for
the complex integrals.
\end{proof}

\begin{proposition}\label{prop:limit}
For $T\geq0$ and real $\eta$,
\begin{align*}
 F(a_n(T)\eta)^n&\longrightarrow e^{-T\eta^4},\\
 e^{n(F(a_n(T)\eta)-1)}&\longrightarrow e^{-T\eta^4},\\
 e^{nT(F(b_n\eta)-1)}&\longrightarrow e^{-T\eta^4}.
\end{align*}
\end{proposition}
\begin{proof}
Multiplication of the convergent cosine and hyperbolic cosine series
gives $F(x)=1-x^4/6+O(x^8)$ as $x\to0$; the degree-two and degree-six
terms cancel. Thus
\[
 n(F(a_n(T)\eta)-1)=-T\eta^4+O(n^{-1}).
\]
For the walk, $F(a_n(T)\eta)>0$ eventually, and
$\log(1+y)=y+O(y^2)$ shows that
$n\log F(a_n(T)\eta)\to-T\eta^4$. Exponentiation proves the first
two limits. Replacing $a_n(T)$ by $b_n$ gives
$n(F(b_n\eta)-1)\to-\eta^4$; multiplication by $T$ proves the third.
For $T=0$, the first two sequences are identically one, as is the third.
\end{proof}

These statements establish only the stated exponential-test limits.
Passing a general entire function through the limit needs a separate
integrability and domination argument. No such unrestricted extension
is used in this paper.

\section{Exact positive moments and compact escape}
\begin{lemma}\label{lem:moments}
For $n\geq0$, real $a$, and $r\geq0$,
\begin{align}
 \E|S_{n,a}|^2&=2na^2,&
 \E|S_{n,a}|^4&=(8n^2-4n)a^4,\label{eq:walkmoments}\\
 \E|C_{r,a}|^2&=2ra^2,&
 \E|C_{r,a}|^4&=(8r^2+4r)a^4.\label{eq:poissonmoments}
\end{align}
\end{lemma}
\begin{proof}
Let $R_n=\sum_{j=1}^n\varepsilon_j$ and
$I_n=\sum_{j=1}^n\delta_j$. These variables are independent, with
$\E R_n^2=\E I_n^2=n$. In the fourth power of $R_n$, only four equal
indices and two equal pairs survive sign averaging. Therefore
\[
 \E R_n^4=n+6\binom n2=3n^2-2n.
\]
Since $|S_{n,a}|^2=a^2(R_n^2+I_n^2)$, expansion gives
\[
 \E|S_{n,a}|^4=a^4\bigl(2(3n^2-2n)+2n^2\bigr),
\]
proving \eqref{eq:walkmoments}. The defining Poisson series gives
$\E N_r=r$ and $\E[N_r(N_r-1)]=r^2$ by shifting the summation index.
Hence $\E N_r^2=r^2+r$. Conditioning \eqref{eq:walkmoments} on
$N_r$ proves \eqref{eq:poissonmoments}. All these nonnegative
moments are finite by the same Poisson moment sums.
\end{proof}

The positive second moment in this lemma must be distinguished from
$\E S_{n,a}^2=0$. The latter is a holomorphic cancellation and is not
a variance in the complex Hilbert-space sense.

\begin{lemma}\label{lem:tail}
If $X\geq0$, $0<m=\E X<\infty$, and $\E X^2\leq A m^2$, then
for $0\leq b\leq m/2$,
\[
 \Pp(X>b)\geq\frac1{4A}.
\]
\end{lemma}
\begin{proof}
Split the expectation at $b$ and apply Cauchy--Schwarz:
\[
 m\leq b+\E[X\mathbf1_{\{X>b\}}]
 \leq b+(\E X^2)^{1/2}\Pp(X>b)^{1/2}.
\]
Thus $m/2\leq \sqrt A m\Pp(X>b)^{1/2}$. Divide by $m$ and square.
\end{proof}

\begin{proof}[Proof of Theorem~\ref{thm:positive}]
The transforms follow from Lemma~\ref{lem:step} and
Proposition~\ref{prop:limit}. At $a=a_n(T)$, the common second moment
of the walk and the rate-$n$ compound law is
\[
 m_n=2\sqrt{6Tn}\longrightarrow\infty.
\]
For the walk the fourth moment is $(48n-24)T\leq2m_n^2$;
for the compound law it is $(48n+24)T\leq3m_n^2$, since $n\geq1$.
Given compact $K$, choose $R>0$ with $K\subset\{|z|\leq R\}$.
Eventually $m_n\geq2R^2$. Apply Lemma~\ref{lem:tail} to $X=|S|^2$
or $|C|^2$, with $b=R^2$ and $A=2$ or $3$ respectively.
The event $\{X>R^2\}$ lies in $K^c$, giving the two constants.

This contradicts tightness with any error smaller than $1/12$.
The same eventual inequalities hold after any $\varphi(k)\to\infty$.
To rule out a weak limit directly, let $\rho$ be a proposed probability
limit. The disks $\{|z|\leq R\}$ increase to $\C$, so choose $R$
with $\rho(\{|z|\leq R\})>1-1/48$. Let $\chi$ be a continuous
function equal to one on that disk, zero outside the disk of radius
$R+1$, and between zero and one. Then
$\int\chi\,d\rho>1-1/48$, while every sufficiently large approximating
law has $\int\chi\leq1-1/12$. This contradicts convergence against
the fixed bounded continuous $\chi$, and proves the assertion for both
families and every divergent reindexing.
\end{proof}

The constants are elementary uniform lower bounds, not claimed optimal.
In particular this proof does not claim that the escaping mass tends to
one; the verified conclusion is already enough to exclude weak limits.

\section{Path-space consequences with an explicit topology}
For $n\geq1$, join the successive partial sums
$a\sum_{j=1}^k\zeta_j$ linearly at times $k/n$ to obtain a continuous
path $\Gamma_{n,a}:[0,1]\to\C$. Give $C([0,1];\C)$ the uniform norm.
Its endpoint is $S_{n,a}$.

\begin{corollary}\label{cor:path}
For $T>0$, the laws of $\Gamma_{n,a_n(T)}$ are not tight in the uniform
path topology. More generally, any family of probability measures on a
topological path space admitting a continuous coordinate map with
coordinate laws $\mu_n^T$ or $\mu_{n,a_n(T)}$ is not tight.
\end{corollary}
\begin{proof}
A continuous image of a compact set is compact. If a family of path laws
were tight, for every $\epsilon>0$ one could choose a compact path set
carrying mass at least $1-\epsilon$ under every member. Its coordinate
image would be a compact subset of $\C$ carrying at least that mass under
every coordinate law. This contradicts Theorem~\ref{thm:positive}.
Endpoint evaluation is continuous in the uniform norm, which proves the
first assertion.
\end{proof}

This implication specifies its topology. The formal Poisson trajectory
corollary uses a continuous coordinate on a countable product space.
It is not silently a statement about every Skorokhod topology: evaluation
at a jump time requires its own continuity hypotheses there. No path-law
limit is constructed from the holomorphic transform alone.

\section{Actual complex tilts and the fixed test}
We prove Theorem~\ref{thm:tilt} through formulas valid at every finite
parameter. Fix $a,\xi\in\R$, $r\geq0$, and put
\[
 h_{r,a,\xi}(z)=e^{-r\psi_a(\xi)}e^{i\xi z},\qquad
 d\nu_{r,a,\xi}=h_{r,a,\xi}\,d\mu_{r,a}.
\]
The normalizing factor is a positive real number because $\psi_a(\xi)$
is real, although $\psi_a(\xi)$ need not be nonnegative.

\begin{proposition}\label{prop:finite}
The density is absolutely integrable, $\nu_{r,a,\xi}(\C)=1$, and
\begin{align}
 \int e^{i\eta z}\,d\nu_{r,a,\xi}(z)
 &=\exp\bigl(r(\psi_a(\eta+\xi)-\psi_a(\xi))\bigr),\label{eq:shift}\\
 |\nu_{r,a,\xi}|(\C)
 &=\exp\bigl(r\cosh(a\xi)(1-\cos(a\xi))\bigr).
 \label{eq:variation}
\end{align}
The exponential test in \eqref{eq:shift} is integrable for the total
variation measure. Furthermore
\begin{equation}\label{eq:phasecancel}
 f_\xi(z)h_{r,a,\xi}(z)=|h_{r,a,\xi}(z)|\qquad(z\in\C),
\end{equation}
so $\int f_\xi\,d\nu_{r,a,\xi}=|\nu_{r,a,\xi}|(\C)$.
\end{proposition}
\begin{proof}
Formula \eqref{eq:compoundnorm} gives density integrability. Multiplying
the density by $e^{i\eta z}$ leaves the same positive normalizer times
$e^{i(\eta+\xi)z}$, whose modulus is integrable by that formula again.
The compound transform \eqref{eq:compound} now proves
\eqref{eq:shift}; taking $\eta=0$ proves mass one.

For completeness, the variation of a finite complex density $h$ relative
to a positive measure $\mu$ is $|h|\mu$. The upper bound follows from
$|\int_A h\,d\mu|\leq\int_A|h|\,d\mu$ on every measurable partition.
For the reverse bound, approximate the measurable phase $h/|h|$ on
$\{h\ne0\}$ uniformly by finitely many phases on a partition of the
unit circle into arcs of arbitrarily small diameter. On each resulting
measurable piece the integral differs from its constant-phase integral
by at most that diameter times $\int|h|$. Summing and letting the diameter
vanish proves equality for every measurable set. In particular,
\begin{align*}
 |\nu_{r,a,\xi}|(\C)
 &=e^{-r(F(a\xi)-1)}\E|e^{i\xi C_{r,a}}|\\
 &=e^{-r(F(a\xi)-1)+r(\cosh(a\xi)-1)}\\
 &=e^{r\cosh(a\xi)(1-\cos(a\xi))}.
\end{align*}
Finally, writing $z=x+iy$, one has
$e^{-i\xi x}e^{i\xi z}=e^{-\xi y}=|e^{i\xi z}|$.
The normalizer is positive real, so the same identity holds for the
density. This is \eqref{eq:phasecancel}. Integrating it proves the last
claim. The test depends on $\xi$ alone, not on $a$ or $r$.
\end{proof}

\begin{proof}[Proof of Theorem~\ref{thm:tilt}]
Apply Proposition~\ref{prop:finite} with $a=b_n$, $r=nT$.
Proposition~\ref{prop:limit} gives
$nT\psi_{b_n}(u)\to-Tu^4$ for every real $u$. Subtraction at
$u=\eta+\xi$ and $u=\xi$ followed by exponentiation proves
\eqref{eq:tiltlimit}.

To prove divergence using a quantitative lower bound, observe that
$|b_n\xi|\to0$. For $|x|\leq\pi$, concavity of sine on $[0,\pi/2]$
gives $\sin(|x|/2)\geq |x|/\pi$. Consequently
\[
 1-\cos x=2\sin^2(x/2)\geq2x^2/\pi^2.
\]
Since $\cosh x\geq1$, eventually
\begin{equation}\label{eq:growth}
 nT\cosh(b_n\xi)(1-\cos(b_n\xi))
 \geq\frac{2T\xi^2}{\pi^2}\,n b_n^2
 =\frac{2T\xi^2\sqrt6}{\pi^2}\sqrt n\longrightarrow\infty.
\end{equation}
The coefficient is positive by $T>0$ and $\xi\ne0$.
Equation \eqref{eq:variation} and the phase cancellation give
\eqref{eq:fixed}. A finite complex measure has a finite integral against
the modulus-one function $f_\xi$. Divergence of this one integral excludes
convergence against all bounded continuous tests. The same lower bound
holds after every divergent reindexing.
\end{proof}

The exclusions $T>0$, $\xi\ne0$ are indispensable. If $T=0$ the law is
$\delta_0$. If $\xi=0$ the density is one and the variation is one,
although positive-law compact escape can still occur at positive $T$.
Nor is the fixed-test conclusion merely a restatement that some sequence
of test functions detects growing variation: a single explicitly given
test works for the entire sequence.

\section{Formal correspondence and limits of the conclusions}
\begingroup\raggedright
All formal statements used here are frozen at verified commit
\code{bab688bb6140cd6316c80680095a914ea0d346db}.
The accompanying machine-readable theorem map records exact declaration
names, locations and source hashes. Human proofs reorganize the finite
sign computations and the weak-limit contradiction; they introduce no
new Lean theorem or hidden probabilistic axiom.

The formal continuous walk paths are defined on all of $\R$, are constant
outside $[0,1]$, and carry the compact-open topology. Restriction to
$[0,1]$ gives exactly the polygonal path above. Corollary~\ref{cor:path}
proves its uniform-topology version directly from the verified endpoint
law; the general continuous-coordinate implication is a classical
deduction, not an additional named formal theorem.

The central modules are \code{ComplexRootHeatTransformLimit},
\code{ComplexRootPoissonNoWeakLimit}, \code{ComplexRootTiltedLaw}
and \code{ComplexTiltVariationDivergence}. The theorem map also identifies
the exact finite transform, moment, path, phase-cancellation and variation
declarations underlying each proof.

The Poisson compact-escape result uses rate $n$ and amplitude $a_n(T)$;
the tilted limit uses rate $nT$ and amplitude $b_n$. These normalizations
give the same analytic heat symbol, but are not the same finite law.
They have been kept separate throughout the statements and proofs.

The baseline's complete build and transitive axiom evidence are recorded
and hash-checked in the accompanying verification receipt. This editorial
work does not rerun that full build. The present paper does not prove a
continuum probability process with biharmonic transition laws, an
unrestricted analytic functional calculus, or a stochastic calculus for
an ordinary weak limit. Analytic functional limits remain meaningful
precisely on their proved test domains. Attribution, independent referee
review and a comprehensive novelty assessment remain open.
\par\endgroup

\begin{thebibliography}{9}
\bibitem{BM}
S. Bonaccorsi and S. Mazzucchi,
\emph{High order heat-type equations and random walks on the complex plane},
arXiv:1402.6140 (2014),
\href{https://arxiv.org/abs/1402.6140}{arxiv.org/abs/1402.6140}.
Primary comparison: Theorem 2, Remark 1, the discussion before Theorem 7,
and Hypothesis 10.
\bibitem{BCM}
S. Bonaccorsi, C. Calcaterra and S. Mazzucchi,
\emph{An It\^o calculus for a class of limit processes arising from random
walks on the complex plane}, arXiv:1603.05437 (2016),
\href{https://arxiv.org/abs/1603.05437}{arxiv.org/abs/1603.05437}.
\end{thebibliography}
\end{document}
